\chapter{Ekstrema związane}
\label{ma2:ch:constrained-extrema}

W rozdziale \ref{ma2:ch:free-extrema-convexity} badaliśmy
ekstrema funkcji na otwartych zbiorach. Teraz dopuszczalne
przyrosty argumentu są ograniczone więzami. Dla więzów
równościowych wykorzystamy twierdzenie o funkcji uwikłanej
oraz opis przestrzeni stycznej; pozwoli to sprowadzić
problem do ekstremum swobodnego na odpowiedniej przestrzeni
Banacha.

\section{Więzy równościowe i ich regularność}

Niech $E,Y$ będą rzeczywistymi przestrzeniami Banacha,
$U\subseteq E$ zbiorem otwartym, $f:U\to\R$ funkcją
celu i $G:U\to Y$ odwzorowaniem więzów. Rozpatrujemy
zbiór dopuszczalny $M=G^{-1}(\{0\})$ w $U$.

\begin{definicja}
\label{ma2:def:constrained-extremum}
Punkt $p\in M$ jest minimum lokalnym funkcji $f$ przy
więzach $G=0$, jeżeli istnieje otoczenie $W$ punktu $p$
w $U$ takie, że $f(x)\ge f(p)$ dla $x\in M\cap W$.
Minimum jest ścisłe, gdy nierówność jest ostra dla
$x\in(M\cap W)\setminus\{p\}$. Maksimum definiujemy
przez odwrócenie nierówności. Zastąpienie $W$ przez
cały zbiór $U$ daje odpowiednie pojęcia globalne.
\end{definicja}

\begin{definicja}
\label{ma2:def:regular-equality}
Jeżeli $G$ jest klasy $C^1$, to więzy równościowe
nazywamy regularnymi w $p\in M$, gdy $DG(p):E\to Y$
jest surjektywny i jego jądro $K=\ker DG(p)$ ma
domknięte dopełnienie $L$ w $E$, czyli
$E=K\oplus L$ jako topologiczna suma prosta.
\end{definicja}

Jądro $K$ jest domknięte jako przeciwobraz zera
przez ciągły operator liniowy. Przy surjektywności
$T=DG(p)$ i skończonym wymiarze $Y$ dopełnialność
wynika automatycznie. Istotnie, wybieramy bazę
$(y_1,\ldots,y_m)$ przestrzeni $Y$ oraz wektory
$e_j\in E$ z $Te_j=y_j$. Odwzorowanie liniowe
$R:Y\to E$ dane przez $R(\sum_j a_jy_j)=\sum_j a_je_j$
jest ciągłe, bo jego dziedzina ma skończony wymiar
(twierdzenie \ref{ma2:thm:finite}). Ponadto $TR=I_Y$.
Operator $RT$ jest ciągłym rzutem na $L=R(Y)$:
$(RT)^2=R(TR)T=RT$, zaś $RTx=0$ wtedy i tylko wtedy,
gdy $Tx=0$. Stąd
$x=(x-RTx)+RTx\in K+L$; przecięcie $K\cap L$
jest zerowe, ponieważ $R$ jest prawostronną
odwrotnością $T$. Obie składowe zależą od $x$ w sposób
ciągły, a $L=\ker(I_E-RT)$ jest domknięte.
Dla $Y=\{0\}$ wystarczy przyjąć $L=\{0\}$.

\begin{stwierdzenie}[Lokalny wykres zbioru dopuszczalnego]
\label{ma2:prop:equality-chart}
Niech $G$ będzie klasy $C^s$, gdzie $s\ge1$ lub
$s=\infty$, i niech więzy będą regularne w $p$.
Dla rozkładu $E=K\oplus L$ istnieje otwarte
otoczenie $A$ zera w $K$ oraz funkcja
$\varphi:A\to L$ klasy $C^s$, taka że
$\varphi(0)=0$, $D\varphi(0)=0$ i odwzorowanie
\[
 \sigma(h)=p+h+\varphi(h)
\]
parametryzuje wszystkie punkty $M$ w pewnym otoczeniu
$p$. Przestrzeń prędkości krzywych klasy $C^1$ w $M$
przechodzących przez $p$ jest równa $K$.
\end{stwierdzenie}
\begin{proof}
Przestrzenie $K,L$ są przestrzeniami Banacha z normami
odziedziczonymi. Ograniczenie $T|_L:L\to Y$ jest
bijektywne: jeśli $Tl=0$, to $l\in K\cap L=\{0\}$;
dla $y\in Y$ wybieramy $x$ z $Tx=y$ i rozkładamy
$x=k+l$, skąd $Tl=y$. Ciągłość odwrotności wynika
z \atlasref{FA-11::cor:bounded-inverse11}{twierdzenia o operatorze odwrotnym}.
Oznaczmy tę odwrotność przez $R$.

Na otwartym zbiorze par $(h,l)$ takich, że $p+h+l\in U$,
rozpatrujemy $H(h,l)=G(p+h+l)$. Jest to odwzorowanie
klasy $C^s$, $H(0,0)=0$,
$D_lH(0,0)=T|_L$, $D_hH(0,0)=T|_K=0$.
Twierdzenie \ref{ma2:thm:implicit-function} daje
otwarte otoczenia $A\subseteq K$, $B\subseteq L$ zera
oraz funkcję $\varphi:A\to B$, z własnością
\[
 H(h,l)=0\ \Longleftrightarrow\ l=\varphi(h)
 \qquad ((h,l)\in A\times B).
\]
Z wzoru na pochodną funkcji uwikłanej mamy
$D\varphi(0)=-R(T|_K)=0$. Ponieważ rozkład jest
topologiczny, $W=p+A+B$ jest otwartym otoczeniem $p$
w $U$, a $\sigma(A)=M\cap W$. Rzut na $K$ zastosowany
do $\sigma(h)-p$ zwraca $h$, więc odwrotność
parametryzacji jest ciągła.

Jeżeli $\gamma:(-\delta,\delta)\to M$ jest klasy $C^1$
i $\gamma(0)=p$, to różniczkując $G\circ\gamma=0$
w zerze dostajemy $T\gamma'(0)=0$. Odwrotnie,
dla każdego $h\in K$ krzywa $\gamma(t)=\sigma(th)$
jest określona dla dostatecznie małych $|t|$ i ma
$\gamma'(0)=h+D\varphi(0)h=h$. To dowodzi obu inkluzji.
\end{proof}

W skończonym wymiarze otrzymujemy w ten sposób opis
$T_pM=\ker DG(p)$ zgodny z
twierdzeniem \ref{ma2:thm:tangent-descriptions}.

\section{Twierdzenie o mnożnikach Lagrange'a}

\begin{twierdzenie}
\label{ma2:thm:lagrange-multipliers}
Niech $f,G$ będą klasy $C^1$, a więzy regularne w $p$.
Jeżeli $p$ jest ekstremum lokalnym funkcji $f$ przy
więzach $G=0$, to istnieje jedyny ciągły funkcjonał
liniowy $\lambda\in\mathcal L(Y,\R)$ taki, że
\[
 Df(p)=\lambda\circ DG(p).
\]
\end{twierdzenie}
\begin{proof}
Stosujemy stwierdzenie \ref{ma2:prop:equality-chart}.
Funkcja $F=f\circ\sigma$ jest klasy $C^1$ na otwartym
otoczeniu zera w przestrzeni Banacha $K$ i ma w zerze
ekstremum lokalne. Z twierdzenia
\ref{ma2:thm:extremum-necessary} wynika $DF(0)=0$.
Ponieważ $D\sigma(0)h=h$ dla $h\in K$, reguła
łańcuchowa daje $Df(p)|_K=0$.

Dla $R=(T|_L)^{-1}:Y\to L\subseteq E$ połóżmy
$\lambda=Df(p)\circ R$. Jest to ciągły funkcjonał
liniowy. Dla $x\in E$ mamy $x-RTx\in K$, ponieważ
$TR=I_Y$. Wobec tego
\[
 Df(p)x=Df(p)(RTx)=\lambda(Tx).
\]
Jeżeli $\lambda_1\circ T=\lambda_2\circ T$, to
dla dowolnego $y\in Y$ stosujemy tę równość do $Ry$,
uzyskując $\lambda_1(y)=\lambda_2(y)$. Dowodzi to
jednoznaczności, a w szczególności niezależności
mnożnika od wybranego dopełnienia.
\end{proof}

Przyjmujemy konwencję funkcji Lagrange'a
\[
 \mathcal{L}_\lambda(x)=f(x)-\lambda(G(x)).
\]
Warunek mnożników ma postać $D\mathcal{L}_\lambda(p)=0$
w całej przestrzeni $E$, nie tylko w $K$. Samo istnienie
takiego mnożnika jest warunkiem koniecznym,
nie dostatecznym.

\begin{wniosek}[Postać współrzędnościowa]
\label{ma2:cor:lagrange-finite}
Niech $E=\R^n$, $Y=\R^m$ i $G=(g_1,\ldots,g_m)$.
Jeżeli gradienty $\nabla g_j(p)$ są liniowo niezależne,
to w każdym ekstremum związanym w $p$ zachodzi
\[
 G(p)=0,\qquad
 \nabla f(p)=\sum_{j=1}^m\lambda_j\nabla g_j(p)
\]
dla jednoznacznie wyznaczonych liczb $\lambda_j\in\R$.
\end{wniosek}
\begin{proof}
Macierz $DG(p)$ ma wiersze $\nabla g_j(p)^{\mathsf T}$.
Ich niezależność daje rząd $m$, a więc surjektywność.
Jądro ma dopełnienie w skończonym wymiarze.
Każdy funkcjonał na $\R^m$ ma postać
$\lambda(y)=\sum_{j=1}^m\lambda_j y_j$, gdzie
$\lambda_j=\lambda(e_j)$ dla standardowych
wektorów bazowych $e_j$. Stosujemy
twierdzenie \ref{ma2:thm:lagrange-multipliers}
i porównujemy współczynniki funkcjonałów względem
iloczynu skalarnego w $\R^n$.
\end{proof}

\begin{przyklad}[Konieczność regularności]
\label{ma2:ex:singular-constraint}
Dla $f(x,y)=x$ oraz $G(x,y)=x^2+y^2$ zbiór
dopuszczalny jest singletonem $\{(0,0)\}$.
Początek jest jednocześnie minimum i maksimum
globalnym na tym zbiorze. Jednak $DG(0,0)=0$,
a $Df(0,0)(h,k)=h$, więc nie istnieje mnożnik
spełniający $Df(0,0)=\lambda DG(0,0)$.
Nie wolno zatem bez sprawdzenia regularności
ograniczać poszukiwania ekstremów do równań
mnożników.
\end{przyklad}

\section{Warunki drugiego rzędu}

\begin{twierdzenie}[Warunek konieczny i warunek dostateczny]
\label{ma2:thm:constrained-second}
Niech $f,G$ będą klasy $C^2$, $p\in M$, więzy
regularne w $p$, $K=\ker DG(p)$. Załóżmy, że
$Df(p)=\lambda DG(p)$ dla pewnego
$\lambda\in\mathcal L(Y,\R)$ i połóżmy
\[
 Q(h)=D^2f(p)[h,h]-\lambda(D^2G(p)[h,h]).
\]
Jeżeli $p$ jest minimum lokalnym związanym, to $Q(h)\ge0$
dla wszystkich $h\in K$. Jeżeli istnieje $c>0$ takie,
że $Q(h)\ge c\|h\|^2$ dla wszystkich $h\in K$, to
$p$ jest ścisłym minimum lokalnym związanym.
Dokładniej, istnieją $r,\kappa>0$ takie, że
\[
 f(x)\ge f(p)+\kappa\|x-p\|^2
 \quad(x\in M\cap B(p,r)).
\]
Dla maksimum odpowiednio odwracamy nierówności.
\end{twierdzenie}
\begin{proof}
Niech $\sigma$ będzie parametryzacją klasy $C^2$
ze stwierdzenia \ref{ma2:prop:equality-chart}.
Na jej dziedzinie $F=f\circ\sigma$ jest równa
$\mathcal{L}_\lambda\circ\sigma$, gdyż $G\circ\sigma=0$.
Ponadto $D\mathcal{L}_\lambda(p)=0$.
Z reguły drugiej pochodnej złożenia
(stwierdzenie \ref{ma2:prop:second-chain}) otrzymujemy
dla $h,k\in K$
\[
 \begin{aligned}
 D^2F(0)[h,k]
 &=
 D^2\mathcal{L}_\lambda(p)[h,k]
 +D\mathcal{L}_\lambda(p)D^2\sigma(0)[h,k]\\
 &=
 D^2f(p)[h,k]-\lambda(D^2G(p)[h,k]).
 \end{aligned}
\]
Jeżeli $p$ jest minimum związanym, to $0$ jest
minimum swobodnym $F$ na otwartym zbiorze w $K$.
Twierdzenie \ref{ma2:thm:extremum-necessary}
daje $D^2F(0)[h,h]\ge0$.

Przy założeniu jednostajnej dodatniej określoności
mamy $DF(0)=0$ i $D^2F(0)[h,h]\ge c\|h\|^2$.
Twierdzenie \ref{ma2:thm:extremum-sufficient}
daje dla dostatecznie małych $h\in K$
\[
 F(h)\ge F(0)+\frac c4\|h\|^2.
\]
Ponieważ $\varphi(0)=0$ i $D\varphi(0)=0$,
z definicji różniczkowalności
$\|\varphi(h)\|/\|h\|\to0$ przy $h\to0$, $h\ne0$.
Zmniejszamy otoczenie tak, aby
$\|\varphi(h)\|\le\|h\|$. Wtedy
$\|\sigma(h)-p\|\le2\|h\|$, skąd
$f(\sigma(h))\ge f(p)+(c/16)\|\sigma(h)-p\|^2$.
Odwrotność $\sigma$ jest ciągła, więc istnieje
$r>0$, dla którego każdy punkt $M\cap B(p,r)$
ma postać $\sigma(h)$ z tego zmniejszonego
otoczenia. Przyjmujemy $\kappa=c/16$.
Gdy $K=\{0\}$, punkt $p$ jest izolowany w $M$
na mocy lokalnego opisu i nierówność zachodzi
bezpośrednio. Dla maksimum stosujemy wynik
do $-f$ i mnożnika $-\lambda$.
\end{proof}

W skończonym wymiarze wystarczy dodatnia określoność
formy $Q$ na $K\setminus\{0\}$, ponieważ jednostajna
dodatnia określoność wynika ze zwartości sfery w $K$
(stwierdzenie \ref{ma2:prop:finite-positive-coercive}).
W przestrzeniach Banacha nie wolno w ogólności
zastąpić oszacowania z jedną stałą $c>0$ samą
nierównością $Q(h)>0$ dla $h\ne0$.

\begin{przyklad}[Ekstrema funkcji liniowej na okręgu]
\label{ma2:ex:lagrange-circle}
Niech $f(x,y)=x+y$ i $G(x,y)=x^2+y^2-1$.
Na $M$ gradient $\nabla G=2(x,y)$ jest niezerowy.
Równania mnożników przyjmują postać
\[
 1=2\lambda x,\qquad 1=2\lambda y,\qquad
 x^2+y^2=1.
\]
Pierwsze równanie wyklucza $\lambda=0$; odejmując
pierwsze dwa równania dostajemy $x=y$.
Stąd $p_\pm=\pm(1,1)/\sqrt2$ i
$\lambda_\pm=\pm1/\sqrt2$.
W punkcie $p_\pm$ jądro $DG(p_\pm)$ jest prostą
$\{(u,v):u+v=0\}$. Dla każdego $h\in\R^2$
\[
 Q_\pm(h)=-2\lambda_\pm\|h\|_2^2.
\]
Twierdzenie \ref{ma2:thm:constrained-second}
daje ścisłe maksimum lokalne w $p_+$ i ścisłe
minimum lokalne w $p_-$. Są one również globalne:
nierówność Cauchy'ego--Schwarza daje
$|x+y|\le\sqrt2$ na $M$, a wartości w $p_\pm$
wynoszą odpowiednio $\pm\sqrt2$. Równość w przypadku
dodatniej (ujemnej) wartości wymaga, aby $(x,y)$
był dodatnią (ujemną) wielokrotnością $(1,1)$,
więc oba ekstrema globalne są jedyne.
\end{przyklad}

\begin{przyklad}[Hesjan funkcji celu nie wystarcza]
\label{ma2:ex:constraint-curvature}
Dla $f(x,y)=y$ i $G(x,y)=y-x^2$ w punkcie $p=(0,0)$
mamy $Df(p)(u,v)=v=DG(p)(u,v)$, zatem $\lambda=1$.
Ponieważ $D^2f(p)=0$, sam Hesjan funkcji celu
nie wykrywa minimum na paraboli. Natomiast
$K=\{(u,0):u\in\R\}$ i
$D^2G(p)[(u,0),(u,0)]=-2u^2$, skąd $Q(u,0)=2u^2$.
Warunek dostateczny jest spełniony. Bezpośrednio
na zbiorze dopuszczalnym $f(x,x^2)=x^2$, co daje
ścisłe minimum globalne w początku.
\end{przyklad}
